\section{Package map and reproduction}
\label{sec:reproduction}

The archive is rooted at
\texttt{unknot\_active\_faces\_batched\_descent\_20261009/}.
The top-level README gives exact commands and the integration contract.
The manifest records source revisions, incoming dependency hashes, and
the identities of retained fixtures. Relative paths are used throughout
the executable research scripts.

\begin{longtable}{p{0.30\textwidth}p{0.63\textwidth}}
\toprule
Path & Contents \\
\midrule
\endhead
\texttt{article/} & Complete modular \LaTeX{} source, final PDF, and vector
performance figure. \\
\texttt{code/fast/} & Pinned supporting implementation, additive modules,
focused and maintained tests, benchmark and replay scripts. \\
\texttt{code/synthesis/data/} & The two pinned historical fixtures needed
by maintained test imports. \\
\texttt{code/reports/} & Historical oracle helpers required by existing
tests, with provenance recorded. \\
\texttt{integration/} & Additive patch against the pinned tree, file
lineage, and the patch-check result. \\
\texttt{evidence/geometry/} & All 24 fixed-batch inputs, source bindings,
expansion records, both witnesses, and raw paired timing samples. \\
\texttt{evidence/counterexample/} & The realized shared-vertex example,
portable construction, retained 24-tetrahedron seed, and batch witness. \\
\texttt{evidence/validation/} & Complete test logs, independent dynamic
audit, cold-census timings, environment, and replay summaries. \\
\texttt{tools/} & Figure regeneration and package verification helpers. \\
\bottomrule
\end{longtable}

From \texttt{code/fast/}, the fast certificate checks do not need an LP
solver or optional graph library:
\begin{lstlisting}
python -S active_research/replay.py
python -S batched_descent_research/replay_evidence.py
\end{lstlisting}
The full test suite is run with:
\begin{lstlisting}
python -m unittest discover -s tests -v
\end{lstlisting}
NetworkX enables the five general matching tests; the source-specific
small-cover tests still run with site packages disabled. Exact geometric
and arithmetic replay use only included modules. The optional regeneration
of the barycentric torus from its original triangulation library is
separate from replay of the retained seed.

The checked producer benchmark can be rerun with the supplied
\path{batched_descent_research/benchmark_batch.py}; it reconstructs
source-bound inputs and performs alternating paired runs. Its command-line
help describes size and output controls. The peeled-census benchmark takes
the retained evidence directory as an explicit argument, and
\texttt{active\_research/benchmark.py} reconstructs its masks and support
queries from the supplied source file. Reproducing exact elapsed times is
not expected across hosts; reproducing exact face pairings, heights,
supports, objective values, and valid certificates is the relevant
correctness condition.

Compile the article from its directory with:
\begin{lstlisting}
latexmk -pdf -interaction=nonstopmode -halt-on-error article.tex
\end{lstlisting}
The final PDF was rendered and visually checked in addition to successful
\LaTeX{} compilation. The package does not require downloading the three
large incoming archives: the few unchanged prerequisite modules are
included and explicitly attributed. It also does not include a new claim
that the complete recognizer has quasipolynomial worst-case complexity.

\section{The general two-endpoint algebraic extension}
\label{sec:general-reward-extension}

The exported algebraic selector allows nonnegative costs $r_i$ and at most
two nonnegative endpoint rewards $c_{vi}$ even when the geometric bound
$c_{vi}\le r_i$ is absent. This extension is not needed for the main
geometric theorem. Here is the correctness argument for its implementation.
Let
\[
 b_v=\max\bigl(0,\max_i(c_{vi}-r_i)\bigr).
\]
A two-endpoint item $i$ on distinct vertices $u,v$ receives adjusted edge
weight $c_{ui}+c_{vi}-r_i-b_u-b_v$. Discard nonpositive edges and retain the
largest parallel edge. Then the optimum of $H$ is the sum of the $b_v$
plus the maximum adjusted matching weight.

To prove the upper bound, take any omitted set and assign each vertex's
positive maximum reward to one attaining item. Delete items assigned no
vertex; their costs are nonnegative. Each remaining item is assigned one
or two vertices. A singly assigned item contributes at most $b_v$. The
doubly assigned items form a matching on their assigned pairs and
contribute their adjusted weights plus those vertices' baselines. Unused
nonnegative baselines can only enlarge the bound. Dropping nonpositive
adjusted edges can also only enlarge it.

For the reverse bound, choose a maximum adjusted matching and, for each
unmatched vertex with positive baseline, choose a baseline-attaining item.
Take the union of these items with the matching items. At each vertex
the actual maximum reward is at least the assigned one; an item chosen
more than once is paid for only once, so its actual total cost is at most
the sum charged in the assignment. Thus $H$ is at least the proposed
matching-plus-baseline value. The upper bound proves equality. In the
geometric case all baselines vanish, and this reduces to the main theorem.
Outside that case the secondary matching-cardinality rule need not minimise
the number of reconstructed singleton items; the API records this limit.

